% ================================================================
%  第 13 章 A) 遗产与继承人 (原书 p.259-280)
%  13.1 身后的学生 | 13.2 孤儿们 | 13.3 时尚——或名人传
%  源文: _work/src2/source.md
% ================================================================
\providecommand{\fnoteplainbody}[3]{%
  \begingroup\renewcommand{\thefootnote}{#1#2}\footnotemark
  \renewcommand{\thefootnote}{#1}\footnotetext{#3}\endgroup
}

\chapter{A)~遗产与继承人}

\section{I.~身后的学生}

\subsection{1.~教学的失败(2)——或创造与自负}

\notehead{44$'$}{[该注由第(I)部分《愚妄与更新》第八章“孤独的历险”第~50~节所召,原书第~227~页]}

\pmark\origpage{173} 这段文字,在读了“一段过去的重量”这最后一节的那位朋友那里“叮”地一声响了\fnote{1}{(*)}{(5月10日)那位朋友不是别人,正是Zoghman~Mebkhout;他慨然允准我揭开我在本注所引那封信(1984年4月2日)的出处上本来觉得应当维持的匿名。}。他写道:“在你许多昔日学生身上,如你所说,那个咄咄逼人、甚至近乎破坏性的‘老板’的形象,依然强烈。于是就有了你的那种印象。”(我想,这里说的“印象”,就是这一节以及补足它的注~n\textsuperscript{o}s~46、47、50 的某些段落里所表达的东西。)他在信的更前面写道:“首先,我认为你暂时离开数学[\,!],离开得好,因为你和学生之间当时有某种隔膜(当然,Deligne除外)。他们都有点被惊呆了……”。

这是第一次,关于我1970年以前“老板”角色的这种论调传到我耳朵里——超出了客套的恭维!还在同一封信的更前面:“……我明白了,你的昔日学生[按:指那些‘1970年以前的’]不太明白数学的\textbf{创造}是怎么回事,而你或许有一份责任……的确,在他们那个年代,问题都已经摆好了……”\fnote{2}{(**)}{(5月10日)前面这段引文被删削得很厉害,是出于对通信人匿名的尊重。完整的引文以及对其真意的评说,见下一条注;起初因缺乏更详尽的信息,我没有领会这段话的真正含义。}。

我的通信人无疑是想说:是我在提“问题”,连同一并要发展的那些观念,而不让我的学生们自己去找到这些问题与观念;也正是在这一点上,我或许遮蔽了他们,使他们不识数学创造工作中那要紧的部分。这恰与一次谈话留下的印象相合\dmark{}\origpage{174}——那是与两位1970年\textbf{之后}的前学生的谈话,前面一条注(注(23iv))里谈到过。的确,我在向我走来的学生身上,首先寻找的是\textbf{合作者},来发展那些在我之内已然成形的直觉与观念,去给一辆已经在那里的大车“推轮子”;他们因而不必从某种虚无中(如我的通信人当年不得不做的那样)去拉动它。然而——让一种可触的、柔软而致密的东西,从不可触之物的迷雾中取得形体——这恰恰从来就是数学工作于我最迷人的方面,也是工作里我尤其感到“创造”在进行的那一部分:某种比一个单纯的“结果”更精微、更要紧的东西的“诞生”。

有时,当我看到昔日学生中的某一位,以轻蔑的态度对待这件弥足珍贵的东西,在他身上抖开 J.H.C.~Whitehead所说的那种“势利”(就是瞧不起自己“本来也会证明”的东西)\fnote{3}{(*)}{见注“年轻人的势利——或纯洁的捍卫者”(n\textsuperscript{o}~27),原书第247页。},我无疑并非全然无关——这样那样地,我脱不了干系。我教学的失败,在1970年之后的一段时期里触目经心,如今又以一种不同而更隐蔽的形状,显现在我第一时期的教学里——尽管按通常的说法,那一期呈现为完完全全的成功!这一点,近年来我不时有所瞥见,也在给几位前学生的信里提到过,而直到如今,没有从他们任何一位那里得到真正的回音。

不过,倘若说我交给学生们做的、他们同我一起做的工作,是纯技术性的、纯粹例行公事式的,无力调动他们的创造才能,这样说怕也不确切。我交到他们手里的,是一些切实而可靠的出发点,他们在其中享有全部的选择自由,可以由此跃起,正如我自己在他们之前曾经做过的那样。我不记得曾给哪位学生出过一个连我自己也不会乐于去做的题目;他们中任何一位与我同行时走过的艰难路段,也都不比我自己在数学家生涯中独自走过的那些更艰难;只要明白工作必须做、也没有别的路,我就没有因此气馁或尥蹶子。

\pmark\origpage{175} 同时我也觉得,我今天所确认的这场失败,其原因比我所出题目的类型、比这些题目在多大程度上含混或分明,更为精微。我在这场失败中的那一份,似乎更应归咎于我对待数学时自身的自负态度;这些态度,我在这场思考中已有机会检视过。即便不曾浸透与某位学生相伴的那项工作本身,它们也必定或多或少浸染了我本人周围的氛围或空气。自负,哪怕以世界上最“不起眼”的方式表达,总是指向一种封闭,一种对事物精微的本质及其美的麻木——无论那是“数学之物”,还是有生命的、我们本可以接待、鼓励、或者也可以从我们高处的优越俯视的人;对我们身边吹拂的气息及其毁伤性的效力——毁伤他人也毁伤我们自己——麻木不仁。

\subsection{2.~一种不义与无力的感觉}

\notehead{44$''$}{[这条注的出现,不遵守写作的时序]}

(5月10日)借着朋友所给的自由引证权——他的信中凡我认为有用的段落都可以引用——我在这里给出一段完整得多的引文\fnote{4}{(*)}{见前一部分“教学的失败(2)——或创造与自负”的第二条脚注,n\textsuperscript{o}~44$'$。},把那段被删削的引文放回它真正的语境:

\begin{quote}
“的确,75到80年间我非常孤立,除了向Verdier提过寥寥几个问题。但那段时间我并不怨你的昔日
学生们,因为没有人真正懂得这条纽带[按:即离散系数与连续系数之间的纽带]的重要性。一切都在
1980年10月改变了:人们发现了这条纽带对半单群的第一个头等重要的应用,即 Kazhdan--Lusztig
重数公式的证明——其中关键地用到了我们所论的范畴等价。这个等价在得到“Riemann--Hilbert对
应”的名字时,没有一句别的评说——毕竟,这再自然不过了!正是在那里,我明白了你的昔日学生
们不很明白数学的\textbf{创造}是什么,而你或许有一份责任。我至今仍有一种不义与无力的感觉。的确,
在他们那个年代,问题都已摆好。这个定理的应用数目,无论在平展拓扑的框架还是超越的框架里,
都很可观,却始终顶着“Riemann--Hilbert对应”的名字!我有种感觉,我的名字配不上这个结果——对\origpage{176}
许多人,尤其对你的昔日学生而言。但正如你从我的工作的引言里能清楚看到的,正是你的“对偶”
形式主义自然地引出了这个结果。不过同你一样,我并不为“可构造离散系数”与晶体系数(或完整
$\mathcal{D}$-模)之间这条纽带的将来担忧。显然,它无论在空间的上同调还是分析中,都适用于许许多
多的领域。”
\end{quote}

正是朋友的信中这一段,启发了我(除本注之外)后来的那道注“当值无名之辈与好上帝定理”。照这封
信的说法,我全然没有料到(如我在适当的地方所解释的),朋友身上那种“不义与无力的感觉”,所针
对的不仅是那种盲目的轻蔑态度——它\textbf{系统性地贬低}他的贡献(这种态度,在我某些昔日学生身
上,后来我也变得颇为熟悉了)——而且是一场真正的诈骗行动:干脆\textbf{变走}一条关键定理的父亲身
份。这一处境,仅仅八天前才向我揭示——关于这一点,见“不义——或归来的意义”及其后各注
(n\textsuperscript{o}s~75至80),它们集结在“讨论会——或Mebkhout层与反常”的标题之下。

\notehead{45}{}

由于我改变了圈子与生活方式,同旧日朋友们会面或其他接触的机会变得稀少。但这并没有妨碍
“拉开距离”的迹象以种种方式表现出来,彼此强弱不一。而在另一些人那里——如Dieudonné、Cartan、
Schwartz,实际上还有所有在我起步时如此善待我的“长者”——我没有感到丝毫这类东西。除他们之
外,我有种感觉:除极少数人之外,我旧日的朋友或学生在数学世界里,其与我之间的关系(无论有无机会表达),在我
退出那个曾经共同的圈子、共同的世界之后,都变得分裂了,带上了“矛盾的双重态度”。

\section{II.~孤儿们}

\subsection{1.~我的孤儿们}

\notehead{46}{[该注由第(I)部分《愚妄与更新》第八章“孤独的历险”第~50~节所召]}

\pmark\origpage{177} 我想借这个机会,在这里说几句关于数学观念与思想的话——在我带到光天化日之下的所有
这些东西当中,它们在我眼里(而且远远领先于其他)具有最广大的分量(注46$_{1}$)\fnote{5}{(*)}{读者可在注~n\textsuperscript{o}~46$_{1}$至46$_{9}$中找到对本注所巡视诸观念的若干更技术性的评说。此外,在我引入的\textbf{诸}观念之外,读者还可在注 n\textsuperscript{o}~88“尸身”里找到关于我著作中(在我“完全做到底”的那一部分著作之内)被其作者视为“执牛耳者”的那一部分的思考。}。首先
是五个紧密相联的关键观念,我将按其特殊性、丰富性(与深度)递增的次序,快速巡视一遍。

首先是同调代数中\textbf{导出范畴}的观念(见注48,原书第274页),以及把它用于一个“通用”的
形式主义——所谓“\textbf{六种运算的形式主义}”(即诸运算 $\overset{L}{\otimes}$、$Lf^*$、$Rf_!$、$R\mathcal{H}om$、
$Rf_*$、$Lf^!$)(注46$_{2}$)——
用于几何学中迄今引入的最重要的几类“空间”的上同调:“代数”空间(如概形、概形叠等)、解析空间
(复解析的,以及刚性的解析的及同类者)、拓扑空间(当然,还在等待各种“温和空间”的语境,肯定还
有其他许多,比如以小范畴的范畴(Cat)作为同伦模型的那一类……)。这一形式主义所涵盖的,既有
离散性质的系数,也有“连续”系数。

这一对偶形式主义及其无所不在性的逐步发现,是通过一场孤独、执着而严格的思考完成的,它从
1956年延续到1963年。正是在这场思考中,导出范畴的观念逐步成形,以及它在同调代数中所担当
的角色的理解也逐步成形。
这套形式主义所缺的,是我的上同调形式主义视野里对“离散系数”与“连续系数”之间那条可
以想见的纽带的理悟——超越局部系统及其以带可积联络的模或晶体模来解释的那个熟知情形之外
的纽带\origpage{178}。这条深刻的纽带,最初在复解析空间的框架里得到表述,是由Zoghman~Mebkhout
(将近二十年后)发现并建立的:以一边由“可构造”离散系数、另一边由“D-模”或“微分算子复形”
观念所构成的导出范畴来表述(见注46$_{3}$)。

将近十年间,由于得不到那些最有资格鼓励他的昔日学生的鼓励——用他们的关切和他们在与我
接触中积累的经验来支持他——Zoghman~Mebkhout在几乎完全的孤立中,坚持他那出色的研究。但这
并没有妨碍他揭示并证明一门正在普遍的冷漠中蹒跚诞生的新晶体理论的\textbf{两把钥匙定理}
\fnote{6}{(*)}{(6月7日)Mebkhout提醒我,在这两个定理之外,还应添上第三个,同样以导出范畴表述,即他(或许
不太恰当地)所谓的 D-模的“\textbf{双对偶定理}”——三者中最难的一个。关于Mebkhout的观念与结果及其
应用的概观,见 Le~Dung~Trang 与 Zoghman~Mebkhout, Introduction to linear differential Systems, Proc. of Symposia in Pure
Mathematics, vol.~40 (1983), part.~2, p.~31--63。}
——两者(这实在是个不好的兆头!)都以导出范畴表述:其一如方才所指,给出“可构造离散系数”与
晶体系数(满足某些“完整性”与“正则性”条件)之间的范畴等价(注48$'$);其二是“\textbf{那}条晶体
的整体对偶定理”,关于从一个光滑复解析空间(不必紧,这意味着额外的可观的技术困难)到一个点
的常值应用。这些是深刻的定理\fnote{7}{(**)}{(5月30日)第二个定理的证明,遇到超越语境惯有的技术困难,
需要求助于所谓“évétesques”技术(原书用语,此处暂保留法文);我猜它可以归入“困难”证明之列。第一个定理的证明是“显然”的——而且深
刻,动用了Heisuke~Hironaka奇点消解的全部力量。如我在“团结”一注(n\textsuperscript{o}~85)倒数第二段所提示的,定理
一旦提炼出来,“随便哪一个”知情的来者都能证明它。可对照“年轻人的势利——或纯洁的捍卫者”
(n\textsuperscript{o}~27)一注中Whitehead的评语。我写那道注的时候,仿佛在某种隐秘的先知的不声不响的口授之下,
还不曾料到现实会在多大的程度上超过我那些怯生生的、摸索式的提示!},为新照亮解析的与概形的
空间(目前为特征零)\origpage{179}的上同调,并许诺这些空间的上同调理论一场规模宏大的更新。它们最终为它们
的作者——在两次进入法国国家科学研究中心的申请被拒之后——换来一个研究员职位(相当于大
学的助理或助教职位)。

这十年间,没有人想到跟Mebkhout说一声“六种运算的形式主义”——我那些学生人人烂熟
\fnote{8}{(*)}{他们是在 SGA~4 与 SGA~5 的讨论班里第一手学到的,或经由 R.~Hartshorne的《留数与对
偶》(Residues and Duality)转手。}、却在任何地方都找不到整理成文文本的形式主义——而他正苦苦对付着超
越语境所带来的巨大的技术困难。他最终是去年从我口中得知它的存在的(以一份看来只有我
一人掌握的公式总表的形式),其时他正不辞辛劳、好意向我讲解他所做的一切;而我,早已不那么在
行上同调了……也没有人想到向他提示:也许先“转向”特征零的概形语境会更“合算”——在那里,
超越语境固有的困难消失,而理论的基本概念问题反而显得愈发清晰。也没有人想到向他指出(或者
哪怕注意到一件我自引入晶体之时就已知道的事
\fnote{9}{(**)}{(5月30日)不过我来得及把它忘了——直到去年与Mebkhout的第二次会面,才凭那次重逢之
力重新想起。(见“冥世重逢”一注,n\textsuperscript{o}~78。)}):光滑空间(解析的或概形的)上的“D-模”,不多
不少,正是“模的晶体”(只要对双方都不问“凝聚”与否的问题)——而后者是一个万能的观念,对带任
何奇点的“空间”与对光滑空间同样行之有效(注46$_{4}$)。

以Mebkhout所展现的才能(与那份少见的勇气),有一点对我说来再清楚不过:倘若处身于一种友善
的氛围,他建立特征零概形的晶体上同调语境中完整的“六种运算”形式主义,将毫无困难,而且会乐
在其中——这样一项宏大计划所需的全部本质观念(包括他自己的,Sato学派与我自己的),在我
看来本已集结在手。以他的才具,这不过是几年功夫的问题,正如平展上同调的一个通用形式主义的
发展是几年功夫的问题(1962--1965)——只要六种运算这条主线(外加换基的两把钥匙定理)已经在
手\origpage{180}。诚然,那是被同侪与见证者的热情与善意的潮流所托举的年月;而不是逆着那些大权
在握者的傲慢自足、逆流而上的劳作……

我来到我想谈的第二对观念:\textbf{概形},以及与它血缘最近的\textbf{topos}。后者是\textbf{site}观念的
更内蕴的版本——site 是我为了使“局部化”的拓扑直觉形式化而最先引入的。(“site”这个词,后来
是Jean~Giraud引入的;他也为使 site 与 topos 的观念具备全部必要的灵活性做了大量工作。)正是代数几
何的种种触目的需要,使我接二连三地引入了概形与 topos。这对观念,蕴含着对代数几何、算术,以及
拓扑的同样宏大的更新之潜能:通过一种\textbf{综合},把这些太久彼此隔绝的“世界”统一在一个共同
的几何直觉里。

以概形的观点、site(或“下降”)的语言,外加十二年关键的奠基工作(还不算我的学生们和其他
志愿者的劳作)对代数几何与算术的更新,二十年来已是既成之事:概形观念、概形的平展上同调观
念(即便还不是平展 topos 与平展叠的观念),终于进入了习惯,进入了公共遗产。

相反,那场同样会把拓扑囊括进来的宏大综合——尽管二十年來,它所需的本质观念与主要
技术工具在我看来已然集结待命\fnote{10}{(*)}{(5月15日)这些“本质观念与主要技术工具”已在 SGA~4
与 SGA~5 两场讨论班的宏大画卷中集结,时在1963至1965年间。这幅画卷的 SGA~5 部分,其成文与出
版所遭遇的离奇变故——十一年之后(1977年)以一种面目全非、满目疮痍的形状问世——为这幅宏
大视野落入“某种时尚”之手的命运,提供了一幅惊心的图像——或者不如说,落入某些我的学生之手:
正是他们最先确立了这种时尚(见下一条脚注)。这些变故及其意义,在最近四个星期的思考中逐步显
露,延续于“伙伴”、“清场”、“格格不入者”、“信号”、“翻转”、“沉默”、“团结”、“蒙骗”、“死者”、“屠杀”
诸注之中(n\textsuperscript{o}s~63$'$'$'$、67、67$'$、68、68$'$与84--88)。}——却还在等待它的时刻。\dmark{}\origpage{181}十五年来(自我告
别数学舞台以来),topos 这个丰饶的统一观念、这个强大的发现工具,被某种时尚\fnoteplainbody{11}{(*)}{(5月13
日)写完这几行之后的六个星期里,思考的继续显明:这种“时尚”最先是由我的某些学生确立的——恰
恰是那些最有条件把某种视野、一些观念与技术手段据为己有的人;他们选择了占有这些劳作的工
具,同时对催生这一切的视野、以及对这视野在其之内诞生的那个人,一并拒认。}打入所谓严肃
观念的冷宫。直到今天,对自己这门科学这种潜在的可观扩展、以及它所提供的新资源,有丝毫觉察的
拓扑学家仍然寥寥。

在这幅更新的视野里,拓扑学家日常摆弄的拓扑空间、微分流形等,连同概形(他听说过)与拓扑、微分或概形的叠(这些却无人谈论),不过是同一类非凡几何对象——\textbf{赋环 topos}(注46$_{5}$)——的诸多化身:拓扑、代数几何与算术
的种种直觉,汇聚于其中的“空间”,共同构成一幅几何视野。各种“模”叠——只要你睁着眼睛,就
会处处遇到(注46$_{6}$)——为它提供了数不清的醒目例子。对它们的深入研究,是深入这些几何对象(或
别的对象,如果有不是几何的对象……)的本质性质的一条头等主线:这些模叠所描述的,正是那些对
象的变动、退化与泛化的样态。然而这份丰富至今无人理会,因为使我们得以精细描述它的观念,不
在通行认可的范畴之列。

这场被拒认的综合\fnote{12}{(**)}{(5月13日)这场综合之遭“拒认”——无论其精神,还是使之成为可能的那
个关键观念——首先出自的,不是别人,恰恰是那项技术手段(概形语言与平展上同调理论的构造)的主
要使用者和受益者:终其全部著作,他受惠于这些手段。这就是Pierre~Deligne。凭他非凡的(源于其非凡
才能的)影响力,以及他相对于我的著作所占据的十分特殊的位置——他像是我的著作的一位隐含
遗产继承人——对我所引入的主要观念(概形观念与平展上同调除外)所设的那道不动声色的系统的
路障,收效极大,无疑在那场\textbf{埋葬}了这些观念的“时尚”的确立中扮演了头等角色;这些观念,已被
迫过了将近十五年的萎顿生活。他的著作,被这层暧昧深深地打上了烙印;这层暧昧,我在接续本注的思
考中第一次瞥见。(见“拒受遗产——或矛盾的代价”,注 n\textsuperscript{o}~47,原书第271页)我离去之后Deligne
著作中这道持续的障碍,这一鲜活而仍然含混的第一知觉,在整个关于\emph{葬礼}的思考中,以惊人的方
式得到了精确与确证;在那场葬礼中,我的朋友扮演着主祭的角色。}还带来一个意外的一面:\dmark{}\origpage{182}某些最
常见空间(注46$_{7}$)的熟知的同伦不变量(或更确切地说,它们的投射有限紧化),被发现配有无人料到的
算术结构,尤其是某些投射有限Galois群的作用……

然而,将近十五年来,在“大世界”里,谁要是贸然说出“topos”这个词——除非是开玩笑,或者他有
逻辑学家的托词——以居高临下的姿态鄙视他,竟成了合乎体面的做法。(逻辑学家是出了名的“与常人不同”,他们的怪癖人们
自会原谅……)以导出范畴的瑜伽来表达拓扑空间的同调与上同调,同样不曾进入拓扑学家中间:对他
们来说,Künneth 公式(对不是域的系数环)至今仍是一个由两条谱序列(或充其量一大串短正合列)
构成的系统,而不是某个合宜范畴中唯一的一个典范同构;换基定理(对真态射,或对光滑态射)他们同
样始终无视——而这些定理(在其近邻平展上同调的框架里)曾构成这门上同调“强力启动”的关键转
折(见注46$_{8}$,原书第270页)。当曾经为发展这门瑜伽出过力的人自己早就把它忘得一干二净,还对那
个想要使用它的可怜人冷眼相待的时候,我就不必惊讶了\fnote{13}{(*)}{(5月13日)后来的思考显明:处境开
始随1981年6月的吕米尼讨论会而改变:在那里,人们看到那些曾经“忘记”(或者不如说,埋葬了……)
这些观念的人,带着它们招摇过市——却依然对那位没有其人便绝不会有那场光辉讨论会的同一“可
怜人”冷眼相待。(关于这场难忘的讨论会,见注 n\textsuperscript{o}s~75与81。)}

第五个观念——也许比其他任何观念都更贴近我的心——是“\textbf{Motif}”。它与前四者不同之处在
于:“\textbf{那}个”好的 Motif 观念(哪怕只在基域上,更不必说任意基底概形),至今尚未成为某个令人
满意的定义的对象,即便为此承认人们会需要的一切“合理的”猜想也无济于事。或者不如说,分明
\dmark{}\textbf{那}条应当首先做出的“合理猜想”,就是\textbf{一条理论之\textbf{存在}}的猜想——它满足这样这样的
数据与性质,而对一个\origpage{183}“在局中”的人\fnote{14}{(*)}{(5月13日)我终于明白:直到今天,唯一(除我
之外)符合这个相当特殊的“多少在局中”的意义的人,是Pierre~Deligne。他有四年之久,一面听“我在
代数几何里所知道的那一点点”,一面日常地做我的 Motif 思考的心腹之交。的确,这些事我跟周围许多
别的同事也谈过,但看来没有一个足够“在行”,能消化那在我之内发展了数年的总体视野,或把
我的提示当作出发点,自己去发展一套视野与纲领(正如我自己当年从Serre的某些思想所给的两三
份“强印象”出发所做的那样)。也许我错了,但我觉得,关心代数簇上同调的人们,只要Deligne——他在
上同调中一言九鼎,而同时又是唯一被认为了解这些 Motif 底细的人——自己还在对 Motif 沉默不表,
就并不处于“把 Motif 当回事”的心理状态。(6月8日)核实之后发现,我最初的 Motif 思考可追溯到六
十年代初——因此它们持续了将近十年。},把它完全显式地陈述出来,本不是什么难事(而且会非常迷
人!)。其实在“离开数学”前不久,我离做这件事只差一步。

从某些方面看,处境与微分积分学英雄时代的“无穷小”颇为相似——不过有两点不同。首先,今天
的我们拥有建造精微数学理论的经验与一套有效的概念装备,那是我们的先辈所缺的。其次,尽管有这
些手段,而这一分明本质的观念出现已二十多年,却没有人屈尊(或敢于,尽管有人不屑于……)
撸起袖子,提炼出 Motif 理论的大致轮廓,正如我们的先辈直截了当地为无穷小计算所做的那样。然而
如今,对 Motif 就像当年对“无穷小”一样清楚:这些生灵是存在的;只要你关心代数簇以及这种簇的族
的上同调,尤其关心它们的“算术”性质,它们就在代数几何中随时显形。也许比对我谈到的那四种观念
更甚,“Motif”——所有观念中最特殊、最丰富的一个——系于一大堆形形色色的直觉:绝非含混,而往
往可以\dmark{}以完美的方式\origpage{184}表述出来(必要时,偶或承认一些 Motif 式的前提)。这些“Motif 式”
直觉中最迷人的一个,在我这里,是“Motif Galois群”:它在某种意义上,使人得以在(绝对意义的)域
与有限型概形的投射有限Galois群上“配一个 Motif 结构”。(为给这一观念以确切意义所需的技术工
作——以那些为 Motif 观念提供暂时基础的“前提”来表述——是在Neantro~Saavedra论“塔纳卡
范畴”的论文里完成的。)

时下的共识,对 Motif 观念比对它的三个受难兄弟(姐妹)(导出范畴、所谓“六种运算”的对偶形式
主义、topos)要稍许客气些:它没有被干脆斥为“涂鸦”\fnote{15}{(*)}{如我在前面一条脚注里提过的,导出范
畴三年前曾享受一场锣鼓喧天的掘墓礼(只是没有提我的名字)。topos 与六种运算还在等待它们的时
刻,Motif 也是——除了两年前被部分开棺出土、换上另一位父亲的那一小块(见注n\textsuperscript{o}s~51、52、59)。(5月13日)},但实际上这毫无区别:只要没办法“定义”一个 Motif、没法“证明”点什么,严肃的人
们就只能绝口不谈(怀着最大的遗憾——这话是听说了的,可人要么严肃要么不严肃……)。当然,只要
人们宣称谈它本身都不严肃,那就永远别指望建造 Motif 理论、就它“证明”点什么!

可那几个在局中的人(以及制造时尚的人)自己心里非常清楚:借助那些保密的前提,可以证明许多
东西。也就是说,从今天起——实际上,自从这一观念在Weil猜想的余波中出现起(它们倒是被Deligne
证明了,这总算是个加分项!)——\textbf{Motif 的瑜伽}就实实在在地存在。但它有着一门\textbf{秘密科学}的
地位:入门者肯定寥寥\fnote{16}{(**)}{(5月13日)我现在似乎终于明白,截至1982年,所谓“极少数入门者”,其实仅有Deligne一个人。的确,他披露了这门“秘密科学”中透过某些重要结果显现出来的部分;这些结果包含在这套瑜伽之中,他随着自己能够证明它们而逐一公布,借此获取功劳,同时隐瞒仍然保密的灵感来源。可是,十五年来竟无人行动起来,最终投身于一门规模宏大的Motif理论,这分明说明,我们的时代远远没有无穷小计算英雄时代那种大胆的活力!}。它尽管“不严肃”,却仍然\dmark{}\origpage{185}使这些罕见的入门者能够在许
许多多上同调的处境中说出“人们有权从中期待什么”。由此衍生出许许多多的直觉与局部猜想;借助“瑜伽”提供的理解之光,人们有时事后便能用手头现成的手段处理它们。Deligne的若干项工作受到这套瑜伽启发\fnote{17}{(*)}{(5月13日)我终于对文献稍有了解。如今我看到,Deligne的全部著作都扎根于这套瑜伽。我的文献抽样以及其他交叉核对使我推测:在他的全部著作中,唯一提及这一来源的地方,是1970年《霍奇理论I》中的一句简短文字,将我与Serre一并提到。(见注n\textsuperscript{o}s~$78'_{1}$、$78'_{2}$。)},其中一项(如果我没有弄错)是他最早发表的工作:证明代数簇的射影光滑态射的Leray谱序列退化(为证明所需,取特征零)。这个结果是由有关“权”的考虑所
提示的;这类考虑具有算术性质——这些正是典型的“Motif 式”考虑,只不过被表达在 Motif 的“几何”的术语里。Deligne证明
这个命题靠的是Lefschetz--Hodge理论,而对那动机(49)(如果我记得不错)只字不提——没有它,谁也不会想
到去怀疑如此难以置信的东西!

Motif 的瑜伽,恰恰最初正是从我得自Serre的这套“权的瑜伽”之中诞生的\fnote{18}{(**)}{我从Serre那里(六十
年代初?)得到的,是一个观念或直觉:有一个重要的东西要理解!它起了一次初始推动的作用;此后几
年的思考紧随其后,先围绕一个“权的瑜伽”,不久又围绕一个 Motif 瑜伽。}。是他让我领略了“Weil猜想”的全部魅力(后来有了“Deligne定理”)。他向我解释了(在假定所论特征下存在奇点消解的
前提下)人们如何能凭权的瑜伽,给任意域上每个概形(不必光滑也不必真)配上“虚拟 Betti 数”
——这一点当时给我极深的震动(注46$_{9}$)。我相信,正是这个观念,是那场围绕权的瑜伽的思考的起
点;那场思考在随后几年里,一直在我奠基成文的任务的边缘上继续。(这也是我在七十年代重新拾起
的思路,借助任意基底概形上的“虚Motif”观念,旨在至少为虚Motif建立“六种运算”的一个形式主义
\origpage{186}。)如果说这些年里我一直对Deligne(我的得天独厚的对话者)以及凡愿听的人\fnote{19}{(*)}{(4月10日)我觉得,Deligne是唯一“听进去”的人——而且他小心地把听进去的东西的专属特权留给了自己。诚然,另一方面,写下这些收尾文字时,我“落在”了事件后面:两年前,Motif瑜伽已有一场部分的开棺出土,却对我的角色不着一字!对此可参见由一项意外发现所引发的注n\textsuperscript{o}~50、51、59——那项发现,对这场持续了十二年的下葬的意味,投下一束(至少对我而言)意外的光。直到那时,我还只是相当混乱地觉察到某种下葬的存在,却不曾费心凑近去看看……}讲这个
Motif 瑜伽,那当然不是为了叫他们把它当作一门只属于他们自己的秘密科学。($\Rightarrow$
注47,原书第271页)

\notehead{46$_{1}$}{}

除了以当年我对Riemann--Roch定理所给的表述(以及我据此给它的两种证明),还有它的各种变体而
引入的观念与观点之外,我几乎不作例外。如果我记得不错,这类变体是在 SGA~5 讨论班1965/66年
的最后一次报告里展开的;而那次报告连同同一讨论班的其他若干报告,都遗失得一干二净。最有意思的一种,我看是可构造离散系数的一条变形;我不知道它此后在文献里是否得到过显豁\fnote{20}{(**)}{(6月6日)我重新找到了它——以一个近旁的形态,顶着“Deligne--Grothendieck猜想”这个动听的名字——在MacPherson1974年发表的一篇文章里。详情见注n\textsuperscript{o}~87$_{1}$。}。值得注意:这一条同样容许一个“Motif”变形,它实质上等于断言:对互异的素数$\chi$(与剩余特征互素),与$\chi$-adic可构造层相配的“示性类”(在正则概形$Y$的Chow环里),当这些层出自同一个“词”——同一个motif(比方说,都是对一个给定的$f:X\to Y$的$R^if_!(\mathbb{Z}_\chi)$)——时,全都相等。

\notehead{46$_{2}$}{}

这个形式主义,可以看作同调中一种“整体对偶”形式主义的某种精粹;以其最“有效”的形
状——撇清了所论“空间”与应用的任何光滑性前提,或态射的真性——它还需要用一个局部对偶形
式主义来补足:在其中人们区分所谓“\textbf{对偶化}”的对象或“复形”\dmark{}\origpage{187}(在 $Lf^{!}$ 运算下稳定
的观念)——并由此有一个“\textbf{双对偶定理}”(以 $R\mathcal{H}om$ 表述),适用于满足合宜的有限
性条件(关于次数,以及局部层面的凝聚性或“可构造性”)的系数。当我说“六种运算的形式主义”时,
我通过后者所暗指的,是对偶性——其“局部”与“整体”两个方面。

深入理解上同调中对偶性的第一步,是六种运算的形式主义在第一个重要情形——诺特概形及上同调凝聚的模复
形——中的逐步展开。第二步是发现(在概形的平展上同调中)这个形式主义对离散系数同样适用。这
两个极端情形,足以确立这种对偶性之\textbf{无处不在}的信念——对Poincaré型“对偶”的一切几何处
境的信念;这一信念,此后又得到Verdier、拉米、吕盖等人的工作的印证。当那股十五年
来压在这套形式主义上、阻碍其大规模使用的\textbf{阻碍}崩解之后,它必将在任意类型的系数上再
度得到确认。

这种无处不在,在我看来是一个分量极重的\textbf{事实}。它使我们不能不感到Poincaré对偶与Serre对偶之间有一种深刻的统一性——Mebkhout最终以所需的一般性确立了它。这种无处不在,使“六种运算的形式主义”成为
同调代数中理解“全方位”上同调对偶现象的基础结构之一\fnote{21}{(*)}{有兴趣的读者,可在本卷的附录里找到这一形式主义的一个概要。}。至于这套颇为精微的结构竟未被行业内
的人显式陈述出来——连“三角范畴”的“好”观念(Verdier的版本仍是一种十分暂时且不充分的形状)
也不曾有人提出——这一事实毫不改变什么;拓扑学家乃至偶或关心上同调的几何学家,仍然竞相无视
这套对偶形式主义的存在,以及奠基它的导出范畴的语言,也同样无济于事。

\notehead{46$_{3}$}{}

D-模与微分算子复形的观点,是由Sato修引入的,最初由他本人及其学派发展,所循的\dmark{}\origpage{188}观点(据我的理解)与Mebkhout所遵循的相当不同;而后者的进路更贴近我的进路。

“可构造性”的诸种观念——对(复解析的、实解析的、分片线性的)“离散”系数——据我记忆,是由我在五十年代末前后首次提炼出来的(几年后我在平展上同调的框架里又重提了它们)。当时我提过一个关于这种观念之稳定性的问题:它对实或复解析空间之间的一个真态射的高阶直像是否稳定——并且不知道这种稳定性在复解析情形是否已被确立\fnote{22}{(*)}{(5月25日)它已由 J.L.~Verdier确立,见“好的参考文献”一注(n\textsuperscript{o}~82)。}。在实解析情形,我当年设想的观念本来就不是对的,因为没有Hironaka的实亚解析集观念——它带有经由直像的本质先行稳定性。至于 $R\mathcal{H}om$ 一类的局部运算,当时清楚的是:在特征零优秀概形的框架里(利用Hironaka奇点消解)确立可构造系数之稳定性的论证,在复解析情形原封不动地成立,双对偶定理亦然(见 SGA~5~I)。分片线性的框架里,自然的稳定性与双对偶定理是“容易的练习”——我当年很高兴做它们,以此检验对偶形式主义的“无处不在”,那正是在平展上同调起步之时(其一大意外恰恰是发现了这种无处不在)。

回到半解析的情形:这一方向上稳定性定理(经由六种运算的可构造系数)的“好”框架,分明是“温和空间”的框架(见《纲领概要》第5、6节)。

\notehead{46$_{4}$}{}

当然,D-模的观点,连同 $\mathcal{D}$ 是一个凝聚环层的事实,为模的晶体揭示了一种比我所惯用的更隐蔽的“凝聚性”观念;它在不必光滑的空间(解析的或概形的)上仍有意义。称之为“\textbf{M-凝聚性}”($M$ 即Mebkhout)才算公道。于是,对一个多少“在局中”(而且健全的数学直觉丝毫不爽)的人来说,以下应当相当显然:那条推广了光滑情形“微分算子复形”的\dmark{}“好的系数范畴”\origpage{189},不应是别的,正应是模的晶体的“M-凝聚”导出范畴(一个晶体复形称为 M-凝聚的,若其各上同调对象如此)。它在没有光滑性假设下仍有合理的意义,并且应当同时涵盖寻常的“连续”(凝聚)系数理论与“可构造”离散系数理论(对后者引入合宜的完整性与正则性假设)。倘若我的视野是对的,Sato--Mebkhout理论相对于先前已知的晶体语境而言,其两个新的概念配料,就是模的晶体的这种 M-凝聚性观念,以及关于 M-凝聚晶体复形的(更深刻的)完整性与正则性条件。这些观念一旦到手,第一桩本质的工作,就是在晶体语境中发展六种运算的形式主义,以便涵盖二十多年前我已发展的两个特殊情形(寻常凝聚与离散)——而我某些前学生上同调学家,早就为了想必更重要的任务把它们忘得一干二净……

Mebkhout其实早已通过研读我的著作,最终得知“晶体”观念的存在;他也感到他的观点(至少在特征零)应当为这个观念提供一条好的进路——但这点建议落入了聋子的耳朵。从心理上说,很难设想他会投身于那桩势在必行的浩大奠基工程:他所处的,恰是那些以上同调权威自居的人的高傲冷漠的气候之中,而他们本是最有资格鼓励——或者说打消……——的人。

\notehead{46$_{5}$}{}

(5月13日)这里说的,主要是以一个\textbf{交换局部环}赋环的 topos。以拓扑空间上这样一个环层的数据来描述“簇”结构的想法,最先由 H.~Cartan引入,又被Serre在其经典工作 FAC(凝聚代数层)中采用。正是这项工作,为引我走向“概形”观念的思考提供了初始推动。Cartan的进路(经Serre采用)所缺的,为涵盖迄今出现的所有类型的“空间”或“簇”所缺的,正是 topos 的观念(即恰好是这样的“东西”:在其上“集合层”的观念有意义,并且具备那些熟悉的性质)。

\notehead{46$_{6}$}{}

\dmark{}作为 topos 的其他出色的例子——它们不是寻常的空间,\origpage{190}而且以“通行”观念似乎也没有令人满意的替代品——我要指出:一个拓扑空间对一局部等价关系所得的商 topos(比如流形的叶状结构,此时商 topos 甚至是一个“多簇”(multiplicité),即局部是一个簇);对几乎任何一种数学结构(至少那些“以有限投射极限与任意归纳极限来表述的”)的“分类”topos。当人们取一种“簇”结构(拓扑的、微分的、实或复解析的、Nash 的,等等,甚至取一个给定基底上光滑的概形结构)时,每种情形都会出现一个特别诱人的 topos,配得上“普遍簇”(所论那一种的)之名。它的同伦不变量(尤其它的上同调——配得上所论簇类的“分类上同调”之名)早就应当被研究和熟知;但到现在,这条路还根本没有被走过……

\notehead{46$_{7}$}{}

这说的是这样一些空间 $X$:其同伦型被“以自然的方式”描述为一个复代数簇的同伦型。后者于是可定义在复数域的一个子域 $K$ 上,$K$ 是素域 $\mathbb{Q}$ 的有限型扩张。投射有限Galois群 $\mathrm{Gal}(\overline{K}/K)$ 于是自然地作用在 $X$ 的投射有限同伦不变量上。往往(比如 $X$ 是一个奇数维同伦球面时)可取 $K$ 为素域 $\mathbb{Q}$。

\notehead{46$_{8}$}{}

(5月13日)当我在Serre的论文 FAC 里学到代数几何的最初几课(正是它把我“触发”向概形)时,换基这个观念本身在代数几何里还几乎无人知晓,除非在换基域这个特殊情形。随着概形语言的引入,这一运算变成了代数几何中最常用的运算,它在每个时刻都会出现。这一运算在拓扑学里至今仍几乎无人知晓(除一些很特殊的情形),在我看来,是拓扑学相对于来自代数几何的思想与技术的孤立的一个典型征兆(在许多其他征兆之中),也是“几何”拓扑学那套不充分基础的一个顽固遗产。

\notehead{46$_{9}$}{}

\dmark{}(6月5日)Serre的想法是:\origpage{191}对一个域 $K$ 上任意有限型概形 $X$,应能配上整数
$$h^i(X)\quad(i\in\mathbb{N})$$
他称之为“虚拟 Betti 数”,使得:
\begingroup\renewcommand{\labelenumi}{\textbf{\alph{enumi})}}
\begin{enumerate}
\item 对 $Y$ 一个闭子概形、$U$ 其补开集,
\[
h^i(X) = h^i(Y) + h^j(U).
\]
{\small [译按:原书此式右端印作 $h^j(U)$,疑应为 $h^i(U)$。]}
\item 对 $X$ 射影光滑,有
\[
h^i(X) = X\text{ 的第 }i\text{ 个 Betti 数}\text{。}
\]
{\small [译按:原书左端印作 $h(X)$;此处按前后定义保留译稿的 $h^i(X)$。]}
\end{enumerate}\endgroup
(比如经由与 $k$ 的特征互素的 $\chi$-进上同调定义)。如果承认 $k$ 上代数概形的奇点消解,那么立刻可见诸 $h^i(X)$ 由这些性质唯一确定。这样一个映射 $X\to(h^i(X))_{i\in\mathbb{N}}$($k$ 固定)的\textbf{存在},利用具真支撑的上同调形式主义,可本质上归约为基域有限的情形。在 $\mathbb{Q}_\chi$ 上有限维的、$\mathrm{Gal}(\overline{k}/k)$ 连续作用的向量空间的“Grothendieck 群”里工作,取 $X$ 的 $\chi$-adic(具真支撑的)Euler--Poincaré特征在该群中的像,$h^i(X)$ 即指 $EP(X,\mathbb{Q}_\chi)$ 的“权 $i$ 分量”的虚拟秩,其中权的观念由Weil猜想导出,外加奇点消解的一个弱形式。即使没有消解,Serre的想法也因Weil猜想的强形式(Deligne在“Conjectures de Weil II”中建立)而得以实现。

我沿着这条路做过启发式的思考,通向“虚拟相对概形”的六种运算形式主义——基域 $k$ 被换成一个多少任意的基底概形 $S$——以及这类($S$ 上有限表现的)虚拟概形的各种“示性类”观念。于是,我(为简化回到基域的情形)被引向考虑比Serre的更细的整数不变量,记作 $h^{p,q}(X)$,满足与上面 a)、b) 类似的性质,并且由通常的公式
$$h^i(X)=\sum_{p+q=i}h^{p,q}(X)$$
给出Serre的虚拟 Betti 数。

\subsection{2.~拒受遗产——或矛盾的代价}

\notehead{47}{[该注直接续注46,见 13.2.1]}


\pmark 值得注意的是,我方才巡视的五种观念之中,有四种(恰恰是那些被当作“不严肃之事”的观念)关系到上同调,\origpage{192}而且首先是\textbf{概形与代数簇的上同调}。不管怎样,这四种观念全都是代数簇的上同调理论的需要向我提示的——先是连续系数,后是离散系数。这就是说,在1955至1970这十五年间,我诸项工作的一个主要动机与一条不变的主导旋律,正是代数簇的上同调。

引人注目的是,这也正是Deligne至今仍视为其灵感主要来源的主题——如果我相信去年IHES小册子就此所说的话\fnote{23}{(*)}{(5月12日)然而,我刚发现,那本小册子里没有任何东西会让读者疑心:我的著作竟与代数簇的上同调、或与任何别的东西的上同调有什么相干!对此可参见即日所写的“葬礼悼词(1)——或恭维”一注(n\textsuperscript{o}~98)。所论小册子,即“有益的挣脱”一注(n\textsuperscript{o}~42)脚注里提到、并在刚提到的“葬礼悼词”注中得到稍细考察的那一本。}。我得知此事时颇为惊讶。诚然,当Deligne(在他关于Ramanujan猜想的漂亮工作(译按:原书人名拼作Ramanuyam,此处按通常拼写订为Ramanujan.)之后)发展了他对Hodge理论的卓越推广时,我还“在场上”,而且完全“在行”。那首先——姑且从这里开始!——对他和对我都一样,是朝着在复数域上以构造性方式建立Motif观念迈出的第一步。在我1970年“转向”之后的最初几年,我自然也听到了Deligne证明Weil猜想的消息(这同时证明了Ramanujan猜想),还有紧随其后的正特征下的“蛮横的”Lefschetz定理。我对他的期待不比这更少!我甚至确信,他应当同时证明了那些\textbf{标准猜想}——那是我在六十年代末提出的,作为(至少)在一域上奠基“半单”Motif观念的第一步,并把这些Motif的若干预期性质翻译成$\chi$-进上同调性质与代数闭链群性质的语言。Deligne后来告诉我,他对Weil猜想的证明肯定无法用来证明那些(更强的)标准猜想,而且他对如何着手也毫无想法。这应当已有十年左右了。从那时起,我再没有听说,在对代数簇上同调之“Motif”(或“算术”)方面\dmark{}\origpage{193}的理解上,发生过其他真正决定性的进展。了解Deligne的手段,我曾默然断定他的主要兴趣必然已转向别处——故而我读到情况并非如此时,深感惊讶。

有一点在我看来是无可怀疑的:至少二十年以来,若不同时在某种程度上扮演“Grothendieck的继承人”,人们已几乎不可能在代数簇上同调的理解上做出大面积革新的工作。Zoghman~Mebkhout已经切身领教了这一点,而(在某种意义上)Carlos~Contou-Carrère亦然——他很快明白,换题目才是上策(注47$_{1}$)。在那些人们无法回避的事情之中,恰恰首推在尽可能贴近Motif的各种系数语境里(这些Motif目前扮演着一种理想“地平线”的角色)发展那著名的“六种运算形式主义”:特征零的晶体系数(沿着Sato学派与Mebkhout的路线,做成“Grothendieck风味”),或特征$p$的晶体系数(主要由Berthelot、Katz、Messing以及一群显然干劲十足的更年轻的研究者研究);Deligne式的“分层pro-模”(promodules stratifiés)(它们显现为凝聚D-模或D-凝聚晶体这一“ind”观念的对偶化变体,即“pro”变体);最后是“Hodge--Deligne”系数(它们看来与Motif一样好,只除了一点:其定义是超越的,且局限于复数域上有限型的基底概形)……在另一端,则摆着这样一个任务:把Motif观念本身从笼罩着它的迷雾中(且不无原因……)提取出来,并且如果可能,去攻克像“标准猜想”这样精确的问题。(关于后者,我曾设想过,除其他途径外,发展一域上射影光滑簇的“中间Jacobi簇”理论,作为或许能得到迹的正性公式的一个办法——那公式是标准猜想的一个本质配料。)

直到我“退出数学”的那一刻,这些任务与问题还在我手里灼烧——滚烫而多汁的东西,其中任何一件、在任何时刻都从不曾在我面前显形为一堵“墙”、一个终止点\fnoteplainbody{24}{(*)}{(5月25日)然而,那本大名鼎鼎的庆典小册子里,一管我自信认得出的匿名手笔,却正是在友好地暗示这一点。参见“葬礼悼词(2)”一注,它接着前一条脚注所引的“葬礼悼词(1)”。}。它们是一个灵感与实质\dmark{}\origpage{194}永不枯竭的源泉——在那里只需朝涌出的地方一拉(而且它“四处涌出”!),如期而至的与不期而至的都会到来。以我这般有限的手段,但只要在工作中不分心,我清楚地知道,只要肯下功夫,一天、一年或十年,人能做成多少事。我也知道——因为我曾在一个他不分心的年代亲眼见过——Deligne的手段是什么,只要他愿意投入,他一天、一周或一个月能做成什么。可是没有人,连Deligne也不能,长此以往地一面从高处俯视那归根结底需要探查的对象本身,连同其语言与整套由某位先行者为这一目的而发展的工具库(何况还是在后者的协助之下,更不必说众多同样参与了劳作的人……),一面做出丰饶的、深刻革新的工作(注59)。

我也想到$M_{g,\nu}$(于$\mathrm{Spec}\,\mathbb{Z}$上)的“Deligne--Mumford”紧化——它参数化亏格$g$、带$\nu$个标记点的连通光滑代数曲线。引入这些对象\fnote{25}{(*)}{载 Pub.~Math.~36, 1969, 第75--110页。评论见注n\textsuperscript{o}~63$_{1}$。},是为了解决在全特征下证明模空间$M_{g,\nu}$的连通性这一问题,办法是从特征零做特化论证。在我看来,这些对象$M_{g,\nu}$(连同群$\mathrm{Sl}(2)$)是我在数学中遇到过的最美、最迷人的对象(注47$_{2}$)。单单它们的存在,带着如此完美的性质,在我眼中就是一桩奇迹(而且是得到了透彻理解的奇迹),其分量远非那桩本待证明的连通性事实可比。对我而言,它们以精髓的方式蕴含了代数几何中最本质的东西,即(几乎所有)一切代数曲线(在一切可以想象的基底域上)的全体——而这些曲线恰恰是一切其他代数簇的终极建筑材料。但所论这类对象——“$\mathrm{Spec}\,\mathbb{Z}$上真且光滑的多重簇”——仍逸出“通行”范畴之外,也就是说,逸出那些人们(出于从不检视的原因)肯于“承认”的范畴之外。常人谈起它们至多闪烁其词,而且一脸抱歉,仿佛还在做什么“泛泛的无稽之谈”,尽管人们小心地会说“stack”或“叠”,以避讳“topos”或“多重簇”这个禁忌的字眼。毫无疑问,这正是这些独一无二的珍宝自十多年前被引入以来,除了我本人在一些未发表的研讨班笔记\dmark{}里之外,再无人(就我所知)研究或利用的原因。取而代之的是,\origpage{195}人们继续要么用“粗糙的”模簇空间,要么用模多重簇的那些有希望是真概形的有限覆盖——然而无论前者还是后者,都不过是它们所自出的那些完美珍宝的某种苍白而跛脚的影子,而后者实际上仍遭放逐……

Deligne的四项工作——关于Ramanujan猜想、关于混合Hodge结构、关于模多重簇的紧化(与Mumford合作)、关于Weil猜想——每一项都是我们对代数簇之认识的一次革新,因而同时也是一个新的起点。这些奠基性的工作相继出现于几年之间(1968--73)。然而将近十年过去,这些巨大的界标并未成为向着预感与未知发起新一轮冲击的跳板,成为更广泛革新的手段。它们通向的是一种阴郁停滞的局面(注47$_{3}$)。这肯定不是因为十年前各家各人手头的那些“手段”像被施了魔法般消失了;也不是因为我们触手可及的那些美好事物突然化为乌有。而是:世界美是不够的——还得有人肯为之欢喜……

\notehead{47$_{1}$}{}

我在这里想到Contou-Carrère在五六年前一个颇有希望的开端:相对局部Jacobi簇的理论,它们与整体Jacobi簇(所谓“广义Jacobi簇”)的联系——对任意概形上的光滑且不必真的曲线概形——以及与Cartier的交换形式群及典型曲线理论的联系。除去Cartier的一次鼓励性回应之外,最有资格赏识Contou-Carrère第一篇注记的人们,对它的接待如此冷淡,以致作者从此绝口不提他储备着的那第二篇,并赶紧换了题目(尽管如此仍未能躲过其他的噩运)\fnote{26}{(*)}{(6月8日)见“棺材3——或有点太相对的Jacobi簇”一注(n\textsuperscript{o}~95)的子注(951)。}。

是我向他建议了局部与整体Jacobi簇这个题目的,作为迈向一个纲领的第一步——那纲领可上溯至五十年代末,尤其指向一个任意维的“阿代尔”对偶复形理论,由局部Jacobi簇构成(对应局部环\dmark{}\origpage{196}任意维的情形),以类比于诺特概形的剩余复形(由其一切局部环的对偶化模构成)。我这一上同调对偶纲领的这一部分(连同其他部分),由于六十年代涌入的种种显得更为紧迫的任务,多少被搁置进了被遗忘的角落。

\notehead{47$_{2}$}{}

确切说来,是所有这些多重簇的族所嵌入的那座“Teichmüller塔”,以及这座塔以基本群胚表述的离散或投射有限范式,构成我在数学中遇到过的最丰富、最迷人的单一对象。群$\mathrm{Sl}(2,\mathbb{Z})$,连同其投射有限紧化的“算术”结构(即$\mathrm{Gal}(\overline{\mathbb{Q}}/\mathbb{Q})$在其上的作用),可视为这座塔的“投射有限版本”的主要建筑材料。译按:原书此处群名的第二个字形似 $\chi$,这里依前文 $\mathrm{Sl}(2)$ 的记号书写。关于此可参见《纲领概要》中的指示(且待《数学反思》诸卷中专论这一主题的卷帙)。

\notehead{47$_{3}$}{}

这一“阴郁停滞”的观察,并不是一个对近十年间围绕概形与代数簇上同调的主要事件了如指掌的人深思熟虑的意见。这是一个“局外人”的单纯的整体\textbf{印象},它是我1982与1983年间同Illusie、Verdier、Mebkhout的谈话与通信中(除其他来源外)得出的。这一印象当然在许许多多方面有待细化。比如,Deligne1980年发表的“Conjectures de Weil II”,就代表一个实质性的新进展,虽说就其主要结果而言算不上一桩意外。看来特征$p>0$的晶体上同调也有进展,更不必说围绕相交上同调的那阵“抢滩”——它终于迫使某些人(尽管不情愿)回到导出范畴的语言,甚至使他们想起那些久遭否认的父系渊源……

\section{III.~时尚——或名人的生活}

\subsection{1.~本能与时尚——或强者的法则}

\notehead{48}{[该注由注46所召,原书第265页]}

众所周知,导出范畴的理论归于 J.L.~Verdier。在他着手\origpage{197}我提议的那项奠基工作之前,我只是以启发的方式、以这些范畴的一个临时定义(后来证明是对的)同它们打交道,对其内部本质结构也只有一个同样临时的直觉(这一直觉在预期的语境里被证明技术上是错的:“映射锥”并不\textbf{函子地}依赖于导出范畴中定义它的那个箭头——它只在相差一个不唯一同构的意义上定义它)。我在五十年代末发展的凝聚层的对偶理论(即凝聚框架里的“六种运算”形式主义)\fnote{27}{(*)}{其中还缺一个对非真态射的$Rf_!$运算(具真支撑的上同调);它是六七年后由Deligne引入的,有赖于他引入的凝聚pro-模语境——在我看来这是一个重要的新观念(后来被他成功地纳入其分层pro-模理论)。},其全部意义唯有在关于导出范畴观念的一项奠基工作之上才得以显豁,而那项工作是Verdier后来完成的。

Verdier学位论文的文本(迟至1967年才答辩),约二十页,在我看来是迄今所写的对导出范畴语言最好的导引,它把这一语言安置在其本质使用的语境之中(其中好几项使用出自Verdier本人)。它只是一项正在撰写的工作的导言,而那项工作最终也是后来才写成的。我可以自诩——即便不是唯一的一个——也是极少数能够作证曾亲手捧读过这部著作的人之一:仅凭这篇导言,其作者被授予国家博士头衔的正当性就得到了确立!这部著作(或者曾经是——我不知道是否还有一份副本存于某处……)是迄今唯一一部按导出范畴的观点给出同调代数的系统奠基的文本。

\pmark 也许只有我一个人在惋惜:无论是那篇导言还是那奠基的正文本身,都未曾\origpage{198}发表\fnote{28}{(*)}{(5月25日)在这些文字写成之后,我发现Verdier学位论文的最初胚芽——写于1963年,即答辩前四年——最终于1967年发表。对此可参见“同伙”与“赊账的学位论文与全险”两注,n\textsuperscript{o}~63$'''$与81。}——以致使用导出范畴语言所需的本质技术行装,散见于文献中三个不同的地方\fnote{29}{(**)}{这三处是:Hartshorne那场著名的凝聚对偶研讨班,其中有我在五十年代下半叶发展的对偶理论至今唯一发表的部分;SGA~4里Deligne的一两个报告;Illusie那部浩繁的学位论文中的一两章。}。这样一个系统参考文本的阙如——其分量本可比肩Cartan--Eilenberg的经典著作——在我眼中,既是1970年我退出数学舞台之后导出范畴形式主义所遭冷落的一个典型\textbf{原因},也是一个典型的\textbf{征兆}。

诚然,早在1968年就已显露(在SGA~5中发展的一种上同调迹理论的需要之际):原初形态的导出范畴观念,以及相应的三角范畴观念,对某些需要而言是不够的,更深彻的奠基工作尚待去做。Illusie在其学位论文中引入“滤导出范畴”,沿这一方向迈出了有用但仍属谨慎的一步(主要为了迹的问题)。1970年我的离去,似乎成了同调代数基础、乃至与之血肉相连的Motif理论基础的全部反思戛然而止、永不再续的信号(注48$_{1}$)。然而就前者而言,为大面积奠基所需的全部本质观念,似乎在我离去之前的年头里就已到手(注48$_{2}$)。(包括“求导者”(dérivateur)这个关键观念,或曰“导出范畴制造机”——它看来是至今所遇到的诸三角范畴之下共同的那个更丰富的对象;这一观念将近二十年后才在一个非加性框架里得到少许发展,载《追寻栈》第二卷的一章。)此外,尚待完成的奠基工作有相当一部分已由Verdier、Hartshorne、Deligne、Illusie做成,这些工作可以原封不动地用于一项综合,把已获的观念纳入“求导者”这个更开阔的视野。

\pmark 诚然,这十五年来\fnote{30}{(*)}{(5月24日)这“逝去的十五年”应当细化——参见注n\textsuperscript{o}~47$_{3}$,以及内容更详的“赊账的学位论文与全险”一注(n\textsuperscript{o}~81)。}对导出范畴观念本身的冷落——在某些人那里,它近于对一段过往的否认——正顺应了某种时尚:\origpage{199}这种时尚对一切奠基的反思都投以轻蔑的目光,无论其多么紧迫\fnote{31}{(**)}{(5月25日)关于这一时尚之出现与延续背后的诸种力量,可参见“掘墓人——或全体教众”一注,n\textsuperscript{o}~97。}。但从另一方面看,我也再清楚不过:平展上同调的发展——今天“人人”都在不加思索地使用它(哪怕只是通过已故Weil猜想的途径……)——若没有导出范畴、六种运算、site与topos的语言(它们最初正是为这一目的而发展)所构成的概念行装,更不必说SGA~1与SGA~2,是绝无可能的。同样再清楚不过的是:今天代数簇上同调理论中可以观察到的那种停滞,若非我昔日某些学生在那些年月里宁肯追随一场他们自己正是始作俑者的时尚、而不肯听从其作为数学家的健全本能——那场时尚久已借着他们的支持获得了法律的效力——原本既不可能出现,更不可能安营扎寨。

\notehead{48$_{1}$}{}

同样的话(带有某些保留)也可用于我那整个代数几何奠基纲领:它只有很小一部分得到实现,而且随我的离去戛然而止。这一停顿给我的震动,尤在于对偶纲领——我视之为特别多汁的部分。Zoghman~Mebkhout的著作,顶风冒雨地继续着,却仍处在这条纲领的脉络里(并因料想不到的观念的汇入而焕然一新)。1976年Carlos~Contou-Carrère的工作(注47$_{1}$,原书第273页已提及)亦然——那是他谨慎地无限期搁置了的工作。此外还有一项关于曲面fppf上同调对偶的工作(Milne)。我所知者,尽于此矣。

诚然,我从未想到要为1955至1970年间为我自行凸显的那项长期工作纲领写一篇概要,一如我为最近十二年写下《纲领概要》。我想原因很简单:从未出现过一个特别的机缘(如同如今我进入CNRS的申请),足以\dmark{}\origpage{200}促成这样一项陈述性的工作。在收入《模型的历史》(《数学反思》II)第I章附录的致Larry~Breen的信(1975年)里,可以找到一些指示,涉及我1970年前日程表上的某些理论(尤其是对偶理论)——这些理论至今仍在等待有人伸出双臂,把它们接引进公共的遗产之中。

\notehead{48$_{2}$}{}

同样的话对Motif理论也成立,只是有一点不同:它多半注定要在一段时期内保持其猜想形态。

\subsection{2.~当值无名之辈与好上帝定理}

\notehead{48$'$}{[该注由注46所召]}

按惯例,人们以完成了提炼并确立一项理论之关键定理的人的名字来命名这些定理;而Zoghman~Mebkhout的名字,看来却被判定配不上这条基本定理——那是四年顽强而孤独的工作(1975--79)的结晶,逆着当时的时尚与其长辈们的轻蔑而行。当这定理的分量再也无法被忽视的那一天,这些人兴致勃勃地把它称作“Riemann--Hilbert定理”,我完全相信(尽管Riemann与Hilbert肯定未曾要求过这份荣耀……)他们这样做自有极好的理由。毕竟(一旦需要之感——即理解一般离散系数与连续系数之间精确关系的需要——顶着普遍的冷漠显现出来,一旦它经由细致而耐心的工作得到提炼与明确,一旦经过相继的几个阶段那正确的陈述终于被提炼出来,一旦它白纸黑字地写下并被证明,而最终,当这件孤独的果实之定理在最不被人期待之处证明了自身——经过了这一切之后)这定理显得如此显然(对那些“本会证明它”的人来说,简直可以径直说“平凡”……),实在犯不着让一个模糊的当值无名氏的名字来占据记忆!

受了这一先例的鼓舞,我提议从今往后把任何一项理论中真正自然、真正基本的定理称作“亚当与夏娃定理”,或者干脆追溯到更远,把荣耀归于当得荣耀之处,径直称之为“\textbf{好上帝定理}”\fnote{32}{(*)}{在我作为数学家的一生中,我从未有过这份乐趣:在一名学生那里激发、哪怕仅仅是鼓励出一篇含有“好上帝定理”的学位论文——至少不曾有过深度与分量可比的。}。

\pmark 据我所知,除我之外,Deligne是Mebkhout之前唯一一个感到如下问题之价值所在的人:\origpage{201}在一个比分层模更宽广的框架里理解离散系数与连续系数的关系,以便能以“连续”的术语诠释任意的“可构造”系数。沿这一方向的第一番尝试,是Deligne1968或1969年在IHES主持的一场(未发表的)研讨班的内容:他在其中引入“分层pro-模”的观点,并给出一个比较定理(在复数域上),把超越的离散上同调与相伴的De~Rham型上同调联系起来——后者对任意特征零基底域上有限型的概形都有意义。(显然,他那时还不知道他那两位遥远的先行者Riemann与Hilbert的卓越结果……)因此,比起Verdier\fnote{33}{(*)}{Verdier作为Zoghman~Mebkhout学位论文的名义导师(并以此身份“恩准过几次讨论”),看来正是(除Mebkhout本人之外)围绕这条基本定理之父系渊源所发生的隐匿——以及在由Mebkhout所发展的D-模观点为代数簇上同调理论开启的革新中其“学生”应得的信誉——的主要当事人。然而我不知道他是否曾为此比Deligne更多一点不安。}或Berthelot\fnote{34}{(**)}{(5月25日)写这些文字时,我(略有踌躇地)没有把我的朋友Luc~Illusie的名字放进这份“本最有资格”向Zoghman~Mebkhout给予那些本属当然的鼓励的我方学生名单里。我当时未曾留意我内心某种不安——它本可以教会我:我这是在暗暗为一位我心怀感情的人开脱,装作卸下他那与我的其他“上同调学家学生”同样背负的责任。},Deligne都处在一个格外有利的位置,足以赏识Mebkhout1975年所投入的研究方向的全部价值,以及此后Mebkhout诸结果——尤其是“好上帝定理”——的价值;这条定理以连续系数的语言,给出了比他本人所提炼的更精妙、更深刻的离散系数领会。但这并未妨碍Mebkhout不得不在一种令人难堪的精神孤立中继续其工作,也未妨碍属于他的信誉(我更愿意说:尤其属于他的信誉)至今——五年之后\fnote{35}{(***)}{(5月25日)其实,这一隐匿首先是Deligne与Verdier本人的作品。对此可参见“不义——或归来的意义”一注,n\textsuperscript{o}~75。}——仍遭藏掖。

\subsection{3.~罐装的权与十二年的秘密}

\notehead{49}{[该注由注46所召]}

\pmark 核对之后(载 Publications~Mathématiques~35, 1968),我发现:在《Lefschetz定理与谱序列的退化判据》一文临近结尾处\origpage{202},有三行文字影射了某些“权方面的考虑”,正是它们使我(以一个略欠一般性的形式)猜想出了该文的主要结果。我怀疑这句谜语式的影射能否对任何人有用,或者在当时能被Serre和我之外的第二个人读懂——反正我们俩本来就已知情\fnote{36}{(*)}{(4月29日)对这篇文章的更细考察——它从不止一个角度看都富有教益——见“驱逐”一注(n\textsuperscript{o}~63)。}。

借机指出:早在六十年代末,在Weil猜想的余波里,一部非常精确的“权的瑜伽”——包括对$R^if_*$与$R^if_!$这类运算下权的行为的刻画——对我说来(因而也对Deligne)已是熟知的了。这部瑜伽的一部分最终在Deligne的“Conjectures de Weil II”(《Publications Mathématiques》1980)中得到确立(在$\ell$-进系数层的语境里;更自然的Motif框架尚待其确立)。如无差错,在前后相隔的这十二年间\fnote{37}{(**)}{(4月19日)我刚收到并饶有兴味地读了Deligne的论著目录,从中发现,早在1974年,他在温哥华大会上的一次报告中就谈到了“权”——如此,“权周围的秘密”是六年,而非十二年。不过在我看来,这一秘密与围绕Motif的类似秘密(1970--1982的十二年间)密不可分。这一秘密的意味,在今天的反思中、在n\textsuperscript{o}~51--52之后那则长篇双注里,刚刚被一束新的光照亮。},文献中没有留下一处踪迹——哪怕是最简短、最不完全的陈述——述说这部(当时还完全是猜想性的)权的瑜伽;在这整段时间里,它始终是寥寥(两三个?)知情人的专属特权\fnote{38}{(***)}{(5月25日)从反思中陆续出现的全部信息来看,这几个“两三个知情人”似乎到头来只有Deligne一人:看来他煞费苦心,把这部得自我的瑜伽的专属收益留给自己,直到1974年(见上一条脚注)——那时机缘已熟,他尽可以把它当作自己地里长出的观念来呈献,既不提我,也不提Serre(见注n\textsuperscript{o}~$78'_{1}$、$78'_{2}$)。

(1985年4月18日)在这些文字写成之后,我又有机会得知Deligne在尼斯国际数学家大会(1970)上的报告“Théorie de Hodge I”(《会议文集》t.1,第425--430页)。与我据手中残缺信息而有理由相信的相反,这篇文章早在1970年就陈述了权的瑜伽的相当一部分。关于这些观念的来源,他只满足于对Serre的一篇文章(它与本题其实无关)以及“Grothendieck的猜想性Motif理论”做了一句谜语式的、纯礼节性的提及。(可与注n\textsuperscript{o}~$78'_{1}$、$78'_{2}$对读。)至于权的观念在$R^if_!$与$R^if_*$这类运算下的行为这一关键问题,则连提都没有提,而且在1980年前引“Conjectures de Weil II”一文之前也不会提;在该文主要定理的关系中我的名字未被提及,同样,在前一条脚注(一年前的今天)所引的那次报告“代数簇上同调中的权”里,Serre的名字和我的名字也都未被提及。}。然而这部瑜伽构成理解代数簇上同调“算术”性质的第一把本质钥匙:它既是在一给定情形中辨认方向的\textbf{手段},也是做出种种预测的手段——其可靠性从未落空——\origpage{203}而同时,正因如此,它也代表着代数簇上同调理论中摆在我们面前最紧迫、最迷人的\textbf{任务}之一。这部瑜伽直到最终被确立(至少在某些重要方面)之前几乎无人知晓,在我看来,正是“\textbf{信息封锁}”的一个格外触目的例证:扮演这一角色的,往往恰恰是那些身居其位、肩负其职、理应确保信息广为传播的人\fnote{39}{(*)}{对此另可参见第32、33节“数学家的伦理”与“注记——或新伦理(1)”,以及与之相关的两注:“共识性的操守与信息控制”与“年轻人的势利,或纯洁的捍卫者”,n\textsuperscript{o}~25、27。}。

\subsection{4.~进步是挡不住的!}

\notehead{50}{[该注由第(I)部分《愚妄与更新》第八章“孤独的历险”第50节所召]}

我在这方面的最初经验,是我为设法让Yves~Ladegaillerie关于曲面同痕定理的学位论文获得出版而付出的那些徒劳努力所结出的意外果实——那项工作,说实话,丝毫不逊于我曾充当“导师”的那十一篇国家博士论文(“1970年以前”的,当然!)。如果我没记错,这些努力持续了整整一年多,主角中有相当一批我昔日的朋友(当然,还少不了我昔日的一个学生)\fnote{40}{(**)}{对此可参见“棺材2——或被切成一段段的剖视”一注,n\textsuperscript{o}~94。}。其主要情节,今天在我眼前仍如一场场通俗笑剧!

这也是我头一回遭遇到某种新精神与新习俗(它们在我昔日朋友圈里已成风气)——在这次反思中我已在这里或那里有机会影射过它们。正是在那一年(即1976年),我头一次(但不是最后一次)获悉:如今,认认真真去证明那些人人都在用、而先行者们从来满足于默认的精微事项(在这里,即曲面拓扑中不存在野性现象)\fnote{41}{(***)}{对此另可参见“注记——或新伦理”那一情节(第33节)。那则大名鼎鼎的“注记”的过错,恰恰在于把一些向来含糊其辞的观念与陈述挑明了;而我却曾一直含蓄地使用它们,来建立那些以我的名字命名的结果,而且将近二十五年来,人人都在毫不脸红地使用这些结果(这一点,那两位鼎鼎大名的同行其实心知肚明)。(6月8日)详情另见“棺材4——或没有鲜花也没有花冠的topos”一注(n\textsuperscript{o}~96)。所谓“以我的名字命名的结果”,是关于某些整体与局部投射有限基本群的生成与有限表现的结果,它们曾在SGA~1等处得到“证明”,所用的下降技术在没有理论上的、细致的论证的情况下仍属启发性质;那样的论证,在Olivier~Leroy那项(显然“无法发表的”)关于topos基本群的Van~Kampen型定理的工作中完成了。},被视为不够严肃(至少就一个无名之辈而言……)。或者,去证明一个把若干已知的深刻定理\dmark{}\origpage{204}作为特例或推论涵盖于其内的结果(这显然表明,那个所谓的新结果只能是已知结果的一个特例或一个容易的推论)。或者,仅仅是在一个结果的陈述里,或是在用一种语言描述另一种情形时,肯费心把自然的假设仔细表述出来(这是可叹的“无谓折腾”的标志),而不肯把自己局限于投某位高格调意见人士之所好的那种特例。(去年我还见到有人责备Contou-Carrère,说他在其学位论文里不该把一般概形、而只把基底域当作立脚点——同时又赏脸承认,他想必是在其临时导师的坚持下才不得不如此。不过,这样说话的人对门道也足够“在行”,本应知道:即便把自身限定在复数域上,证明的需要也会逼着人引入一般的基底概形……)

当今某种时尚的迷途,竟至于不仅鄙薄细致的证明(乃至干脆鄙薄证明本身),甚至常常鄙薄成形的陈述与定义。照纸张与被填鸭的读者的耐心目前的行情,不久之后,人们大概就不再肯为如此昂贵的奢侈品劳神了!顺着当前的趋势外推,应当能够预言这样一个时刻:届时出版物中将不再有人费神去明确陈述定义或定理,人们从此满足于用一些代码词来指称它们,把填补空白的差事留给自己那不知疲倦、天才横溢的读者,随他按自己的悟性行事。审稿人的差事也将随之大为轻松:他只需到“名人录”(Who is Who)年鉴里查一查,作者是否以可信著称(反正谁也没法反驳一个由空白与省略号组成的精彩论文),抑或相反,是一个不便启齿的无名氏——那他早就(如同今天已然如此、且由来已久)被当然地扫地出门……
